% main.tex — audit of the Sony ILCE-7SM3 (a7S III) firmware update file.
% Authors: Oxunjon Ubaydullayev and Aiodam (autonomous research agent).
% Every number below is bound to an artifact by paper_check.py: it is re-derived from the
% artifact bytes, must appear in this source, and must appear in the printed PDF.
% Appendix A is generated from numbers.json by make_numbers.py, not typed by hand.
\documentclass[12pt]{article}
\usepackage[a4paper,margin=0.92in]{geometry}
\usepackage{fontspec}
\usepackage{polyglossia}
\setmainlanguage{english}
% Type: a palatino-style face (TeX Gyre Pagella) with its matching math font, instead of the
% wide, heavy DejaVu Serif the first edition used. Mono stays DejaVu Sans Mono: the paper is full
% of hex dumps and long hashes, and this face has the full ASCII range at a legible weight.
\setmainfont{TeX Gyre Pagella}
\setsansfont{TeX Gyre Heros}[Scale=0.96]
\setmonofont{DejaVu Sans Mono}[Scale=0.84]
\usepackage{unicode-math}
\setmathfont{TeX Gyre Pagella Math}
\usepackage{booktabs}
\usepackage{array}
\usepackage{xcolor}
\usepackage{longtable}
\usepackage[colorlinks=true,linkcolor=blue,urlcolor=blue]{hyperref}
\usepackage{microtype}
\emergencystretch=1.6em
\setlength{\parskip}{0.25em}

\newcommand{\code}[1]{\texttt{\small #1}}
\newcommand{\lc}[1]{\texttt{\small #1}}
\newcommand{\lcf}[1]{\texttt{\footnotesize #1}}
\newcommand{\hx}[1]{\texttt{\footnotesize #1}}
\newcommand{\hl}[1]{\texttt{\footnotesize #1}}



\title{\bf What is openly readable in a Sony ILCE-7SM3 (a7S III)\\
firmware update file, and where the trust boundary really sits\\[0.4em]
\large A static audit of a single file: measurements, limits, and the errors of the
instruments that took them}
\author{Oxunjon Ubaydullayev\thanks{%
Commissioned the audit and wrote its terms of reference; performed the measurements on the
live device (reading the camera's host key and scanning his own network) and supplied the raw
logs; decides what is published.}%
\and
Aiodam (autonomous research agent)\thanks{%
The measurements, the instruments, the format analysis and this text are its work. The errors
of the instruments are not deleted from the record; they are listed in \S\ref{sec:errors}.}%
}
\date{4 October 2026}

\begin{document}
\maketitle

\begin{abstract}
\noindent
We take apart a single public Sony firmware update file, \code{BODYDATA.DAT}, for the camera
ILCE-7SM3 (a7S III): \textbf{743\,818\,632} bytes, sha256 \hx{dbdd8b02\ldots c9eb}. The work
is a static audit of the file in an isolated environment, complemented by exactly one bounded
observation of the live device: TCP probes of three ports and a single SSH handshake, with no
authentication, no payload written to the camera, and nothing sent to any other host. The
result has two halves, and both were measured rather than deduced.

\textbf{First: nothing inside the file is protected by anything the file carries.} The
container is checked by a CRC32 (\hx{0x1fafdcaa}) that any reader can recompute without a key,
and so is the integrity of the image contents. The firmware header is encrypted with a shared
\emph{public} key in ECB mode, so it is readable by anyone holding the published key; and the
mode itself leaks structure: across all 710 MiB of ciphertext there is exactly \textbf{1} run
of identical 16-byte blocks (\textbf{26} blocks, starting at byte \textbf{96}), and that run
was predicted from the decrypted header \emph{before} it was measured. The body
(99.99993\,\% of the file) is decrypted with a \textbf{public} key published in an open but
unmerged pull request, yielding a stream of \textbf{740\,912\,128} bytes, sha256
\hx{11880291\ldots 1e66}, which contains the updater's own file-system image
(squashfs, \textbf{58} files), a \textbf{159}-member firmware tar, and \textbf{1}
unencrypted OpenSSH private key (ECDSA P-256) shipped inside a publicly distributed image.

\textbf{Second: no signature inside the file is verified by anything the file carries.}
Public RSA keys are present --- \textbf{6} distinct moduli --- but none of them verifies the
256-byte block at the tail ($m[0{:}2] = \mathtt{49\,7f}$: no PKCS\#1 v1.5 structure), and no
verification key is present at all. The updater's own pipeline map
(\code{config/config.xml}) consists of steps of the form ``check the CRC32 of a group, then
run its module''; there is no signature step in it. The only honest conclusion is that the
\emph{trust boundary lives outside the file, in the device}, and cannot be probed from a file.

One result came from the live device. The camera presents a host key that matches the key
inside the file in none of three comparison routes (they differ in \textbf{64} of the
\textbf{66} bytes of the point), while negotiating \emph{exactly} the algorithm set that the
firmware's own \code{sshd\_config} imposes. So the firmware's sshd is indeed running on the
camera, and its host key is its own --- which means the private key shipped inside the public
file is \emph{not} this camera's live identity, although it remains a private key in an open
distribution. What could not be established is whether that identity is stable at all: whether
the live key survives a power cycle was \emph{not} measured, and neither was whether it is the
same on every camera of the model. If it is created afresh on every boot, a client cannot pin
the camera and the cryptography is present while the authentication is not.

What this work does \emph{not} establish is as much a part of the result as what it does:
signature verification on the device (the verification key is not in the file), rollback
protection, the parser's robustness against malformed input, behaviour on an interrupted
update, and the debug interfaces of a release build all need a camera, a bootloader dump or
service mode --- not a file. The trust boundary sits outside the file.

Finally, the paper is checked by an instrument rather than taken on trust. Every number is
re-derived from its artifact's bytes (\code{paper\_check.py}); a full re-verification of all
findings by a separate instrument yields \textbf{47} green checks out of \textbf{47} with
\textbf{5} controls that must fail, and finds \textbf{4} corrections to numbers recorded in
earlier reports.
\end{abstract}

\tableofcontents

\section{Subject and scope: what was checked, and what it cannot be}
\label{sec:setup}

There is exactly one subject: a firmware update file \code{BODYDATA.DAT},
743\,818\,632 bytes (709.4 MiB), sha256 \hx{dbdd8b02\ldots c9eb}, md5
\hx{ae4d074d\ldots baac0}, sha1 \hx{03c83a6b\ldots 36b2f}. No experiment modified the
original: copies were made in \code{/tmp}, synthetic images were built separately.

The terms of reference (\code{TZ\_SONY\_ADUDIT.md}) ask, among other things, for the
signature-verification mechanism on the device, the parser's robustness against malformed
input (fuzzing under emulation), rollback protection, behaviour when an update is interrupted,
and whether debug interfaces are present in a release build. From a single file these items
are \textbf{not checkable}: signature verification is done by code in the camera, and error
handling by its bootloader. Everything below is marked as measured on the file; everything
that concerns device behaviour is marked as not measured.

Two limits were observed without exception, and they are stated here because one of them is
not what the word "static" alone would suggest. First, the analysis \emph{of the file} is
static, and it is complemented by exactly one bounded observation of the live device: TCP
connects to three ports and a single SSH handshake, performed on the owner's own camera, with
no authentication, no payload written to the camera, and nothing sent to any other host
(\S\ref{sec:live} reports it in full, including the fact that this observation yielded a
negative result). Second, neither the
paper nor the instruments shipped with it contain \textbf{any key material} --- neither the
body key's value nor the private key. The instruments read the body key from the audit source
at run time and print only \emph{fingerprints}; the paper names fingerprints and the location
where the key was published, never the key.

\section{Method: the rules this work followed}
\label{sec:method}

Four rules shaped the work, and they explain the shape of the paper.

\begin{enumerate}
\item \textbf{Every claim gets a route independent of the one that produced it.} The container
format was parsed with our own code and compared against an independent public implementation
(\code{fwtool.py}); the body was unpacked by two implementations (our byte-wise one and the
library's) and the bytes compared; the key fingerprint was computed by our own P-256
arithmetic and checked with an outside \code{ssh-keygen}; the container CRC32 was recomputed
by our own code.
\item \textbf{Every check has a control that must go red.} An instrument that cannot say
``no'' does not prove ``yes'': a control is the same kind of check, phrased so that the
correct answer to it is \emph{failure}.
\item \textbf{Every number is bound to an artifact.} It is re-derived from the artifact's
bytes on every check of the paper (\S\ref{sec:verification}) instead of being quoted from a
report.
\item \textbf{The errors of the instruments stay in the record.} This work produced
\textbf{27} of them (F140--F166), none deleted; the ones that changed a \emph{conclusion}
rather than the wording are in \S\ref{sec:errors}.
\end{enumerate}

\section{The container: parsed to the byte, but not authenticated}
\label{sec:container}

The file is a sequence of blocks of the form ``\code{u32} length, 4-byte tag, payload''
(the first field is the length of the \emph{payload}, so a block's body starts eight bytes
after the block begins; Table~\ref{tab:container}). The parse agrees with an independent
implementation: the tags \code{DATV}, \code{PROV}, \code{UDID}, \code{FDAT} and \code{DEND}
sit where they should, and these blocks account for the whole file with nothing left over.

\begin{table}[htbp]\centering
\small
\caption{Container blocks (the offset is the block's start, i.e. its length field).}
\label{tab:container}
\begin{tabular}{lrll}
\toprule
tag & offset & payload length & what it is \\
\midrule
\code{DATV} & 8 & 4 & container version \\
\code{PROV} & 20 & 4 & provenance \\
\code{UDID} & 32 & 60 & USB descriptor table \\
\code{FDAT} & 100 & 743\,818\,512 & encrypted body (its body starts at 108) \\
\code{DEND} & 743\,818\,620 & 4 & CRC32 of everything before it \\
\bottomrule
\end{tabular}
\end{table}

The container's outer integrity is a CRC32: the stored value \hx{0x1fafdcaa} equals the
\code{crc32} of the first 743\,818\,620 bytes. That catches accidental corruption, but it is
not authentication: a modified file with a recomputed CRC32 is accepted by the reference
parser, while the same file without the recomputation is rejected (controls C8/C8c in the
instrument \code{checks.py}).

The \code{UDID} table holds \textbf{7} entries of 8 bytes each (PID, VID, flags): six ordinary
descriptors \code{054c:0448}, \code{054c:047d}, \code{054c:0d15}, \code{054c:0d16},
\code{054c:0d19}, \code{054c:0d1a} with flag 1, and a seventh with flag 2 whose PID is
\code{03e2} (\code{054c:03e2}). What matters here is that the update-mode descriptor sits
\emph{inside the file itself} rather than being imported from an outside table: previously it
was known only from the independent \code{Sony-PMCA-RE} project, and now it is confirmed from
within.

\section{The header is readable with a public key, and ECB leaks structure}
\label{sec:header}

The first 512 bytes of the body (that is, starting at offset 108) decrypt under the public
key \code{key\_aes} --- the one that has lain in the open in a public implementation of the
format for years --- in \textbf{ECB} mode. Inside is a coherent structure: a 4-byte frame
(\hx{c4a9fc03}), the magic \code{UDTRFIRM}, format version \code{0100}, mode \code{U}, flag
\code{N}, model \textbf{\code{0x91030083}}, firmware version \textbf{5.01}, region 0, and the
component layout (Table~\ref{tab:header}).

\begin{table}[htbp]\centering
\small
\caption{The \code{FDAT} header: measured fields.}
\label{tab:header}
\begin{tabular}{lr}
\toprule
field & value \\
\midrule
header checksum (\code{crc32} over $[12,508)$) & \hx{0x9eb4163d} \\
bytes $H[508{:}512]$ that this checksum requires & \hx{00 00 00 00} \\
file system \code{U} & offset 512, size 1\,253\,376 \\
slot \code{P} & offset 512, size 0 (no data) \\
firmware area & offset 1\,253\,888, size 739\,658\,240 \\
number of file systems & 2 \\
\bottomrule
\end{tabular}
\end{table}

Two things in this section deserve their own paragraph.

\textbf{(a) The header's integrity is not key-based either.} A \code{crc32} over four trailing
bytes is a bijection, so the stored checksum \emph{determines} the bytes $H[508{:}512]$, and
they are \hx{00000000}. A synthetic header that declares \code{firmware.offset} beyond the end
of the file is silently accepted by the reference parser (control: with a broken CRC, it is
rejected). The format and the reference implementation therefore carry no invariant that
``offsets lie inside the image'' (CWE-20/CWE-1284) --- measured on the implementation, not on
the camera's parser, which we never saw.

\textbf{(b) ECB leaks structure, and this is a prediction rather than hindsight.} From the
decrypted header one can see where a run of identical 16-byte ciphertext blocks \emph{must}
appear: \textbf{26} of them. Measuring the file itself, \emph{without any key}, finds exactly
\textbf{1} such run, starting at block \textbf{6} --- that is byte \textbf{96} of the
ciphertext --- with length \textbf{26} blocks (\textbf{416} bytes) out of the
743\,818\,240 bytes of ciphertext. This is the ECB leak (CWE-327): an observer who does not
know the key still learns the fact ``416 bytes here are identical text''.

\section{The body: opened with a public key; the ``unaccounted bytes'' were arithmetic}
\label{sec:body}

The body key --- two components, \code{key\_aes} plus the 32-byte generation key
\code{key\_cxd90057\_k8} for CXD90057 ---
sits in a public but \emph{unmerged} pull request against a public implementation of the
format (author DavidBuchanan314, opened 10 September 2026). Confidentiality of the body is
therefore broken publicly, and that is not our doing: we merely found a publication that was
not in the main branch. Unpacking yields a stream of \textbf{740\,912\,128} bytes, sha256
\hx{11880291\ldots 1e66}, reproduced independently by two different unpackers (our byte-wise
one and the library's) with identical bytes.

The stream is a sequence of 1024-byte frames: a 2-byte checksum, 2 bytes of size and
end-of-stream flag, then the payload. Measured: \textbf{726\,385} frames, all checksums
valid (\textbf{0} bad), payload size \textbf{1020} bytes in every frame but the last
(\textbf{448}), the end flag set for \textbf{1} frame and only the last, and a tail padding
of \textbf{572} bytes, all \code{0xff}.

This closes a question that the first report carried as ``2\,906\,112 unaccounted bytes''
(hypothesis F-07): they are the frame overhead plus the padding,
$4 \times 726\,385 + 572 = 2\,906\,112$. The mistake was not in the file but in my
coordinates: the header's offsets are in \emph{stream} coordinates (after the frames are
stripped), whereas the ``remainder'' was computed as the difference between the ciphertext
length and the sum of the components. The stream equals
$512 + 1\,253\,376 + 739\,658\,240 = 740\,912\,128$ bytes, leaving
\textbf{0} extra bytes.

\section{What is inside}
\label{sec:inside}

\subsection{The updater's file-system image}

The first 1\,253\,376 bytes of the stream are a \textbf{squashfs} 4.0 image (ZLIB, built on
29 March 2024), sha256 \hx{62e86fad\ldots d8ccaa}. It was unpacked in full: \textbf{58}
files, \textbf{3\,151\,747} bytes, of which \textbf{31} are ELF (all AArch64) and
\textbf{21} are scripts. Roles were read from the files themselves rather than assigned by
name: \code{edisx.elf} partitions the storage, \code{pformat.elf} (\textbf{1\,631\,008}
bytes) writes images and carries a full OpenSSL, including \code{AES\_encrypt},
\code{SM4\_encrypt}, \lc{Camellia\_cbc\_encrypt},
\lc{d2i\_PKCS\allowbreak 8\_PRIV\_K\allowbreak EY\_INFO} and \code{CMS\_EncryptedData\_it}, with its own
\code{[Encrypt] Error:} branch --- image encryption on the device exists, and lives there.
\code{config/chassis} reads \code{CXD90057+CXD90058} (independent confirmation of the
hardware class) and \lc{config/b\allowbreak ody\_vers\allowbreak ion} reads \code{700.104.039}.

\subsection{The update pipeline is CRC32 and nothing else}

\code{config/config.xml} in this image is the update map: every step is tied to its own
checksum and the order is precisely \emph{check the CRC32 of a group, then run the module}.
There is no signature step in the pipeline. The internal \code{.sum} files are CRC32 too (we
recomputed them over the members' own bytes; one changed byte yields a consistent new
checksum without any key). Next to it sits \code{MODE\_SERVICE=2}: a service mode exists, and
in it \code{us\_mkfs.sh} skips repartitioning and writes only the application image. And
separately: in \code{loader\_writer.sh} the md5 check of the loader is
\textbf{commented out}, while \lc{nor\_load\allowbreak er\_write\allowbreak r.sh} declares \code{loader\_nor.md5}
and never calls the check.

\subsection{The firmware archive}

The firmware area (739\,658\,240 bytes) is a \code{tar} of \textbf{159} members: partition
configurations (\code{mount.conf}, built 29 June 2023 --- that is the date of the
\emph{image build at Sony}, not of our work), partition images (WBI, \code{/usr} = \code{/dev/nflasha15}), regional
files, i2c firmware, CA certificates inside the \code{/usr} image, and libraries' string
tables. A note on method: \textbf{172} occurrences of the armour line \code{CERTIFICATE} in
the stream are a count of \emph{literals}, not of objects, and \S\ref{sec:key} shows where
that can lead: of the \textbf{8} blocks that begin with \code{BEGIN PUBLIC KEY}, one turned
out to be a string with no body. The file names its own model: \code{ILCE-7SM3} occurs
\textbf{61} times (including the pair \code{ILCE-7SM3 v5.01}), and the only other model name,
\code{ILCE-9}, occurs \textbf{2} times in a neighbouring table. The public tables that map a
model code to a name do not contain this model: our code \code{0x91030083} is a different
CXD90057-class device from the a7 IV (\code{0x01030081}).

\subsection{One private key inside a public distribution}
\label{sec:key}

In the whole stream there are \textbf{5} lines beginning \code{BEGIN OPENSSH PRIVATE KEY} and
the same number of \code{END} lines, \textbf{10} literals in total. Exactly one of them is a
real key: the other four pairs are string tables in the libcrypto/libssh2 libraries, where
\code{BEGIN} and \code{END} stand tens of bytes apart. The distinguishing measurement is the
body length: a real key has \textbf{435} base64 characters between the lines.

The key sits at offset \textbf{587\,628\,092} inside the member
\code{0700\_part\_image/dev/nflasha5} (the \code{/wbi} partition image), in a block of
\textbf{504} bytes. Parsing it gives
\code{openssh-key-v1}, \code{cipher=none}, \code{kdf=none} (that is, it is \emph{not}
password-protected), type \lc{ecdsa-sh\allowbreak a2-nistp\allowbreak 256}, curve \code{nistp256}, comment
\code{root@(none)}. That it works was checked along three independent paths: our own P-256
arithmetic reproduces the public point from the private scalar; an outside \code{ssh-keygen}
prints the same fingerprint \hx{SHA256:J8L9aBLE\ldots NODI}; and the sha256 of the public blob
is \hx{27c2fd68\ldots c4d3832}. Control: flipping one bit of the scalar breaks $dG = Q$.

What the key is for can be read from the same file. The \code{sshd\_config} member
(\code{0800\_appli/tmp/ssh/sshd\_config}, \textbf{3094} bytes, at offset
\textbf{740\,733\,952}) contains a \emph{single} \code{HostKey} line,
\code{HostKey /tmp\_network/ssh/ssh\_host\_ecdsa\_key}, and exactly four algorithm lines:
\code{HostKeyAlgorithms} \lc{ecdsa-sh\allowbreak a2-nistp\allowbreak 256}, \code{KexAlgorithms}
\lc{ecdh-sha\allowbreak 2-nistp2\allowbreak 56}, \code{MACs} \code{hmac-sha2-256}, \code{Ciphers}
\code{aes128-ctr}. Next to it we found the full text of the script that creates such a key ---
\code{/usr/bin/create\_host\_key.sh}: \code{ssh-keygen -q -b 256 -t ecdsa -N '' -f
/tmp\_network/ssh/ssh\_host\_ecdsa\_key}. The comment \code{root@(none)} is exactly what
\code{ssh-keygen} writes when run as root on a machine with no hostname.

The key grants no access: key-based login is \textbf{forbidden} in that config
(\lc{PubkeyAu\allowbreak thentica\allowbreak tion no}, \lc{PermitRo\allowbreak otLogin \allowbreak no},
\lc{Password\allowbreak Authentication \allowbreak no}, \code{MaxSessions 0}), and only forwards to \code{localhost:\allowbreak15740} and \code{localhost:\allowbreak60152}
are allowed --- it is a server identity for a tunnel, not a way in. Its risk runs the other
way: the private half of a host key sits in a file every camera owner downloads, and those are
the same bytes in every copy of firmware 5.01.

\subsection{Tail and signature: the verification key is not in the file}

After the ciphertext area (the tail is read at offset \textbf{743\,818\,348} inside \code{FDAT})
there remain \textbf{272} bytes: 16 bytes of plaintext IV
(\hl{89a85276c720\allowbreak 8ea533a0e187\allowbreak fac4092c}) and \textbf{256} bytes --- exactly the size of an
RSA-2048 signature. The file does contain RSA public keys (\textbf{6} distinct moduli: two of
3071 bits, three of 3072, one of 2048), and the 2048-bit one (at offset
\textbf{429\,029\,376}) is the only one that could, by size, sign those 256 bytes. It does not:
raising the tail to the power $e$ modulo $n$ produces no PKCS\#1 v1.5 structure (the first two
bytes are \hx{49 7f}), and no PSS either. An exhaustive sweep (2 keys $\times$ 10 candidate
messages $\times$ 4 hashes $\times$ \{v1.5, PSS\}) yields not one acceptance. The published
Sony module from the \code{Sony-PMCA-RE} project fares no better.

The right conclusion is not ``there is no signature'' but \textbf{``there is nothing in the
file with which to confirm one''}: no verification key is present, so it must live in the
bootloader or in ROM. That is the trust boundary, and it lies outside the subject of this
audit.

\section{The live device: the shipped key is not used, and that changes the conclusion}
\label{sec:live}

The decisive check for \S\ref{sec:key} is a single one: read the host key of a real camera and
compare. It was done --- from two machines independently (the owner's Mac and this container),
against the device \code{192.168.1.102}, whose port \code{22} answered in \textbf{22} probes out
of \textbf{22} over three minutes, while port \code{15740} (PTP/IP) was closed.

The result is \textbf{VERDICT: MISMATCH}. The live key is
\hl{SHA256:0VOxJ\allowbreak qn75lys4Lf91\allowbreak tSlZu8nazWWx\allowbreak D6rTVc35jZo8\allowbreak 7U} (blob sha256
\hx{d153b126\ldots 3668f3b5}), against the firmware key
\hx{SHA256:J8L9aBLE\ldots NODI}. They differ in \textbf{64} of the \textbf{66} bytes of the
point: a different key, not a typo. Three comparison routes were used (the blob's base64, the
\code{ssh-keygen} fingerprint, the sha256 of the decoded blob), and all three said ``no''.

But the camera said more than ``my key is different''. In the same SSH exchange it negotiated
\emph{exactly} four algorithms, and that is precisely the set imposed by the \code{sshd\_config}
inside the firmware (\S\ref{sec:key}). An ordinary machine does not choose such a set: in the
owner's run, his own workstation on the same network (\code{192.168.1.103}, OpenSSH 10.3)
negotiated \lc{mlkem768\allowbreak x25519-s\allowbreak ha256} and \code{chacha20-poly1305}, while two other devices
with port 22 open turned out to be \code{dropbear} and negotiated no matching kex at all. Hence
a conclusion worth more than the verdict itself: \textbf{the sshd running on the camera is the
firmware's own, but its host key is its own too}. So the path
\code{create\_host\_key.sh} really does execute on the device, and the private key inside the
public file is \emph{not} this camera's live identity.

What that changes and what it does not, stated plainly. It does not change: the key in the file
remains a private key in an open distribution, and the leak does not disappear because the
device does not use it. It does change: the question moves from ``whose key is this'' to ``is
the camera's identity verifiable at all'' --- and if the key is created afresh under
\code{/tmp\_network} on every boot (which is what the recovered script and the path in
\code{sshd\_config} say, and which we did \emph{not} measure), then the answer is no. A client
that honestly records the fingerprint
in \code{known\_hosts} will see a key-change warning every time the camera is switched on: the
cryptography is present and the authentication is not, and the worst effect here is not
cryptographic but human --- it trains the operator to dismiss the warning.

What we did \emph{not} measure, and do not claim: whether the live key changes across a reboot
(it lives outside persistent storage, but that is a premise read from a config, not a
measurement --- it needs one run before and one after a power cycle), and whether it is the
same on every camera of the model. And separately, as a fact without an explanation: the PTP/IP
port \code{15740}, which a camera normally opens, was closed in that state --- so ``the camera
was found'' rests here on the algorithm set and on the key match, not on the PTP/IP port.

\section{Classification of the findings, and what to do}
\label{sec:findings}

\begin{table}[htbp]\centering
\small
\caption{Findings: what was measured, its class, and what follows.}
\label{tab:findings}
\begin{tabular}{p{0.13\textwidth}p{0.50\textwidth}p{0.30\textwidth}}
\toprule
class & finding & access type \\
\midrule
CWE-327 & a shared public header key in ECB mode; the block run leaks (26 blocks, predicted
ahead of measurement) & local delivery of the file \\
CWE-345 & container and inner integrity are CRC32, recomputable without a key &
local delivery of the file \\
CWE-20/1284 & the format does not require \code{offset}/\code{size} to lie inside the image;
the reference parser accepts an offset past the end of the file & local delivery of the file \\
CWE-321/798 & an unencrypted OpenSSH private key (ECDSA P-256) inside a publicly distributed
image; the same bytes in every copy of 5.01. Measured \emph{not} to be the camera's live
identity (\S\ref{sec:live}) & hygiene; no vector demonstrated: the camera does not use this
key, so impersonation would need a client that pinned it from elsewhere \\
CWE-1188 (\emph{hypothesis}) & if the host key is created afresh under \code{/tmp} on every
boot, the camera's identity cannot be pinned by a client. Whether it survives a reboot was
\emph{not} measured & hypothesis; the decisive measurement is one run before and one after a
power cycle \\
--- & no signature step in the update pipeline; no verification key in the file &
needs the device \\
\bottomrule
\end{tabular}
\end{table}

The recommendations follow from the findings rather than from general principle.

\begin{enumerate}
\item \textbf{Retire the shared header key.} One key across several camera generations means a
single publication opens every header. Replace ECB with an authenticated mode (AES-GCM/SIV)
keyed per model, so that integrity cannot be recomputed without the key.
\item \textbf{Base version decisions only on signed data.} A header that anyone can
manufacture is no basis for deciding whether to install or to roll back: either include the
version inside the signed area, or keep it out of the header entirely.
\item \textbf{Do not use CRC32 as a trust boundary}, and separately, \textbf{enforce bounds}:
today the format does not require offsets to lie inside the image.
\item \textbf{Do not ship a private host key.} Generate it on the device in persistent
storage; and do not keep it under \code{/tmp} if a client is ever to pin the camera's
identity.
\item \textbf{Anti-rollback in practice} (a monotonic counter in a protected area) is not
checkable from a file, but structurally it is absent from the header.
\item \textbf{The next step needs a device, not a file}: a bootloader dump or service mode.
Only with them do the remaining items of the terms of reference stop being ``not checkable
from a single file''.
\end{enumerate}

\section{What was not established}
\label{sec:unknown}

The list is kept without softening, because it is exactly what separates a report from
advertising.

\begin{enumerate}
\item Signature verification on the camera: what is signed, by what, and whether it is checked
before the image is applied.
\item The root of trust (bootloader/ROM): it is not in the file.
\item Rollback protection in practice; behaviour when an update is interrupted.
\item The robustness of the \emph{camera's} parser against malformed input: what we measured
is a property of the format and of the \emph{reference} implementation, not of the camera.
\item What the 256-byte tail formally is (its shape matches an RSA-2048 signature; there is
nothing with which to confirm it).
\item Whether the live host key changes on reboot, and whether it is identical across units
of the model.
\item The authenticity of the file itself: nothing inside it proves it came from Sony.
\item The reason port \code{15740} was closed in the measured state.
\end{enumerate}

\section{How all of this was checked}
\label{sec:verification}

\textbf{(1) The audit instruments, run again.} Every existing check was re-run today against
the unchanged file:

\begin{table}[htbp]\centering
\small
\caption{Re-run of the audit instruments (all against the file with sha256
\hx{dbdd8b02\ldots c9eb}).}
\label{tab:runs}
\begin{tabular}{lrll}
\toprule
instrument & checks & controls & outcome \\
\midrule
\code{audit/checks.py} & 15/15 & 3/3 red & all consistent \\
\lc{audit/re\allowbreak port\_che\allowbreak ck.py} & 2/2 & 3/3 red & all consistent \\
\lc{audit/ro\allowbreak und277\_c\allowbreak heck.py} & 13/13 & 5/5 red & all consistent \\
\lc{audit/ro\allowbreak und278\_c\allowbreak heck.py} & 17/17 & 5/5 red & all consistent \\
\lc{audit/ro\allowbreak und279\_c\allowbreak heck.py} & 28/29 & 2/2 red & one SKIP: the camera is offline now \\
\lc{audit/pa\allowbreak rser\_bou\allowbreak nds.py} & --- & --- & format bounds, as described \\
 & 47/47 & 5/5 red & all consistent \\
\bottomrule
\end{tabular}
\end{table}

The \code{round279\_check.py} row is a measurement, not a caveat: the live half of that suite
(C22--C25: the MISMATCH verdict, the algorithm set, the key matching the owner's key) did
\emph{not} execute today, because the camera does not answer from this container; the
instrument prints SKIP rather than a green. The live measurements were taken earlier and
travel with the report and with this paper as separate witness files.

\textbf{(2) All findings re-verified by a separate instrument.}
 is not a retelling of the reports: it takes 47 statements and
re-derives each one from raw bytes with its own code (its own container parse, its own frame
unpacking, its own \code{tar} walk, its own \code{openssh-key-v1} parse, its own P-256
arithmetic, its own \code{crc32}). Outcome: \textbf{47} non-control checks green out of
\textbf{47}, and \textbf{5} controls red as they must be. The controls are: searching for a
string that is definitely absent; a corrupted header prefix; a corrupted key scalar; a
corrupted frame; a read past the end of the stream.

\textbf{(3) What the re-verification itself found: 4 corrections to our own reports.}
This is the most important part of the section, because they concern numbers already shipped
in delivered packages.

\begin{itemize}
\item \textbf{The key block is 504 bytes, not 501}, and its body is \textbf{435} base64
characters, not 471 (471 is the offset at which the closing line begins, read as a length).
\item \textbf{The offset of \code{sshd\_config} is 740\,733\,952, not 740\,736\,525} (the
recorded value lands 2573 bytes \emph{inside} the file; its size, 3094, was right).
\item \textbf{There are 8 blocks beginning \code{BEGIN PUBLIC KEY}, not 4}, and \textbf{6}
distinct RSA moduli, not 5 (two of 3071 bits, three of 3072, one of 2048). One of the eight
blocks is a literal with no body --- a string table, not a key; the same is true of the single
literal \lc{BEGIN EN\allowbreak CRYPTED \allowbreak PRIVATE \allowbreak KEY}.
\item \textbf{The body of \code{FDAT} starts at byte 108} (the length field is at 100).
Checking such an offset matters: a check that reads the header at byte 100 produces exactly the
same \emph{appearance} of a result (decryption ``works'', hashes come out) on different bytes.
\end{itemize}

None of the corrections changed a conclusion; all four changed numbers that had been recorded.
That is the price of the rule ``every number is bound to an artifact'': while a number lives in
prose, nobody re-measures it.

\textbf{(4) The numbers of this paper.} \code{paper\_check.py} re-derives every
number in the paper from the artifact's bytes (the path and sha256 of each artifact are
recorded in \code{numbers.json}), requires the number to appear both in this source and in the
printed PDF, and \code{selftest\_paper.py} proves the instrument can go red: on a
corrupted number in the text, on a substituted value in a copy of an artifact, and on a deleted
artifact. Appendix~\ref{sec:appendix} is printed from \code{numbers.json} rather than typed, so
the list of numbers cannot drift away from the check. A separate layout instrument
(\code{layout\_check.py}) checks the printed pages for letters lying over other
letters and for ink outside the type area --- a defect TeX reports only as an ``Overfull
hbox'', and not for every case. A third instrument, \code{string\_check.py}, requires the
\textbf{46} names these claims rest on to be present both in the source and in the text
extracted from the printed PDF --- it exists because names, unlike numbers, had gone missing
from the paper without any check noticing (\S\ref{sec:errors}).

\section{The errors of the instruments in this work}
\label{sec:errors}

The work produced \textbf{27} of them; here are those that changed a conclusion, or that
belong to a class worth recognising. The full list is in \code{FAILURES.md} and in the
per-round reports.

\begin{itemize}
\item \textbf{An instrument counted instructions as artefacts.} The claim ``the image has
three \lc{OPENSSH \allowbreak PRIVATE \allowbreak KEY} blocks'' was a count of \emph{literals} inside libraries'
string tables. The real key is one. Same class: ``six public crypters'' --- the seventh lay in
an unmerged pull request, invisible to a search of the main branch.
\item \textbf{An instrument measuring a state it was itself changing.} The first version of
the CRC-tampering check computed the checksum while the write had not yet been flushed to
disk, and returned a false ``the CRC cannot be forged''.
\item \textbf{Precision of comparison.} The report checker compared numbers with a relative
tolerance of $10^{-6}$ and therefore missed \code{739658240} $\rightarrow$ \code{739658241}.
A comparison must be as precise as the text claims to be.
\item \textbf{One false conclusion, reached twice.} Parsing the line
\lc{kex: hos\allowbreak t key al\allowbreak gorithm:\allowbreak (no mat\allowbreak ch)} yielded the word \code{match)}, and
\code{ssh-keyscan} terminates its lines with a carriage return; as a result the check twice
declared ``a different algorithm set'' about a camera that had negotiated exactly the
firmware's set. It was caught only by running against the live device instead of against my
model of a server.
\item \textbf{An instrument that reddened by itself.} The control for the expected frequency
of magic bytes compared three-byte patterns against \code{/dev/urandom} and failed in one run
out of three for reasons unrelated to the subject: a control must be deterministic.
\item \textbf{Lazy decryption.} The ``wrong key'' control was blind, because the library's
\code{decrypt()} is lazy and reads no bytes: it ``accepted'' a key of all zeros.
\item \textbf{A check that looked at the wrong representation.} While building the number
instrument of this paper, the payload offset was verified by comparing the file's raw bytes
against the plaintext magic --- but that region is encrypted, so the check could only ever
fail. The offset is now verified on the \emph{decrypted} window.
\end{itemize}

\section{Reproduction}
\label{sec:repro}

\begin{enumerate}
\item Verify the subject: \lc{sha256su\allowbreak m BODYDA\allowbreak TA.DAT} must give
\hx{dbdd8b02\ldots c9eb}. No audit command writes to that file.
\item Audit instruments: \lc{audit/che\allowbreak cks.py},
\lc{audit/re\allowbreak port\_che\allowbreak ck.py}, \lc{audit/ro\allowbreak und277\_c\allowbreak heck.py},
\lc{audit/ro\allowbreak und278\_c\allowbreak heck.py}, \lc{audit/verify\_find\allowbreak ings.py}; the live ones are
\lc{audit/ro\allowbreak und279\_c\allowbreak heck.py} (needs \code{paramiko}) and
\lc{MAC\_FINGERPRINT.sh} on a Mac.
\item Paper: \lc{publication/run\_all.sh} --- it runs the instruments, rebuilds the
numbers, builds the PDF, binds every number to its artifact, checks the layout, runs the
self-tests and writes the manifest.
\end{enumerate}

\appendix
\section{The numbers of this paper and their artifacts}
\label{sec:appendix}

The table is printed from \code{numbers.json} by the generator \code{make\_numbers.py}: it
\emph{is} the list of numbers that \code{paper\_check.py} checks. The ``artifact'' column is
the path whose bytes the number is re-derived from; the sha256 of every artifact is recorded
in \code{numbers.json}.

\begingroup\scriptsize
\begin{longtable}{p{0.46\textwidth}p{0.50\textwidth}}
\caption{The numbers of the paper, grouped by the artifact whose bytes each one is re-derived from. The sha256 prefix of every artifact is printed once, at the group it belongs to; the full values are in \code{numbers.json}.}\label{tab:numbers}\\
\toprule
number & what it is \\
\midrule
\endfirsthead
\toprule
number (continued) & what it is \\
\midrule
\endhead
\midrule
\multicolumn{2}{l}{\lcf{BODYDATA.DAT}\quad\textbf{sha256 dbdd8b02ed81d6f0}} \\
\midrule
\lcf{743 818 632} & size of the update file, bytes \\
\lcf{dbdd8b02ed81\allowbreak d6f0e0a0d36b\allowbreak 55f5d12fa1bc\allowbreak f8682fbad139\allowbreak 3996f10d1494\allowbreak c9eb} & sha256 of the object \\
\lcf{ae4d074dca26\allowbreak 8917cc969771\allowbreak 4f6baac0} & md5 of the object \\
\lcf{03c83a6b616c\allowbreak 9f2164954373\allowbreak 59fac585db33\allowbreak 6b2f} & sha1 of the object \\
\midrule
\multicolumn{2}{l}{\lcf{-}\quad\textbf{sha256 -}} \\
\midrule
\lcf{709.4} & the same size in MiB \\
\midrule
\multicolumn{2}{l}{\lcf{BODYDATA.DAT}\quad\textbf{sha256 dbdd8b02ed81d6f0}} \\
\midrule
\lcf{8} & DATV block: offset \\
\lcf{20} & PROV block \\
\lcf{32} & UDID block \\
\lcf{100} & FDAT block \\
\lcf{743 818 512} & FDAT payload, bytes \\
\lcf{743 818 620} & DEND block: offset (= length of the prefix the CRC covers) \\
\lcf{0x1fafdcaa} & stored container CRC32 \\
\lcf{7} & entries in the UDID table \\
\lcf{0x03e2} & last descriptor: PID of the updater mode \\
\lcf{108} & offset of the FDAT payload (= start of the body) \\
\lcf{c4a9fc03} & frame of the first body block \\
\lcf{UDTRFIRM} & body magic \\
\lcf{0100} & format version \\
\lcf{0x9eb4163d} & header CRC32 \\
\lcf{5.01} & firmware version in the header \\
\lcf{0x91030083} & model code in the header \\
\lcf{0} & region \\
\lcf{00000000} & the bytes H[508:512] that the CRC32 requires \\
\lcf{512} & offset of file system U in the stream \\
\lcf{1253376} & size of file system U \\
\lcf{0} & size of slot P (empty) \\
\lcf{1253888} & offset of the firmware area \\
\lcf{739658240} & size of the firmware area \\
\lcf{2} & number of file systems in the header \\
\lcf{1} & runs of identical 16-byte ciphertext blocks (measured without the key) \\
\lcf{6} & start of the run, 16-byte block index \\
\lcf{96} & start of the run, bytes from the start of the ciphertext \\
\lcf{26} & length of the run, 16-byte blocks \\
\lcf{416} & length of the run, bytes \\
\lcf{743818240} & length of the ciphertext, bytes \\
\midrule
\multicolumn{2}{l}{\lcf{out/stream.b\allowbreak in}\quad\textbf{sha256 11880291ae4bd08b}} \\
\midrule
\lcf{740912128} & unpacked stream, bytes \\
\lcf{11880291ae4b\allowbreak d08b61c9d084\allowbreak 75ee8a33703b\allowbreak 1c8f08ad35cb\allowbreak cda3a4717410\allowbreak 1e66} & sha256 of the unpacked stream \\
\midrule
\multicolumn{2}{l}{\lcf{-}\quad\textbf{sha256 -}} \\
\midrule
\lcf{726385} & number of stream frames \\
\lcf{0} & frames with a bad checksum \\
\lcf{1020} & payload size of a frame \\
\lcf{448} & payload size of the last frame \\
\lcf{1} & frames carrying the end flag \\
\lcf{572} & padding of the last frame, bytes \\
\lcf{2906112} & frame overhead plus padding (= the former "unexplained" quantity) \\
\lcf{0} & bytes beyond the declared components \\
\midrule
\multicolumn{2}{l}{\lcf{out/fs\_user.\allowbreak sqsh}\quad\textbf{sha256 62e86fad175e0326}} \\
\midrule
\lcf{62e86fad175e\allowbreak 0326969d5f1c\allowbreak ea29d780feb4\allowbreak 96ce8034acd2\allowbreak c9c3954320d8\allowbreak ccaa} & sha256 of the squashfs image \\
\midrule
\multicolumn{2}{l}{\lcf{out/fs\_root}\quad\textbf{sha256 tree:5cb85a8b2b2}} \\
\midrule
\lcf{58} & files in the unpacked updater image \\
\lcf{3151747} & total size of the unpacked files, bytes \\
\lcf{31} & of them ELF \\
\lcf{21} & of them scripts \\
\midrule
\multicolumn{2}{l}{\lcf{out/fs\_root/\allowbreak bin/pformat.\allowbreak elf}\quad\textbf{sha256 3ce2a76f477e9373}} \\
\midrule
\lcf{1631008} & size of pformat.elf, bytes \\
\midrule
\multicolumn{2}{l}{\lcf{-}\quad\textbf{sha256 -}} \\
\midrule
\lcf{159} & members in the firmware tar \\
\midrule
\multicolumn{2}{l}{\lcf{out/stream.b\allowbreak in}\quad\textbf{sha256 11880291ae4bd08b}} \\
\midrule
\lcf{61} & occurrences of the string ILCE-7SM3 in the stream \\
\lcf{2} & occurrences of the string ILCE-9 in the stream \\
\lcf{172} & occurrences of the CERTIFICATE armor literal (literal count) \\
\lcf{587628092} & offset of the private-key block in the stream \\
\lcf{504} & size of the key block (BEGIN..END) \\
\lcf{435} & base64 characters in the key body \\
\lcf{ecdsa-sha2-n\allowbreak istp256} & key type \\
\lcf{none/none} & cipher and KDF (none) \\
\lcf{root@(none)} & comment inside the key \\
\lcf{SHA256:J8L9a\allowbreak BLEimTdDW5m1\allowbreak iV7wU64LP1iW\allowbreak o3+lj9MGCxNO\allowbreak DI} & key fingerprint (ssh-keygen) \\
\lcf{27c2fd6812c4\allowbreak 8a64dd0d6e66\allowbreak d6257bc14eb8\allowbreak 2cfd625a8dfe\allowbreak 963f4c182c4d\allowbreak 3832} & sha256 of the key public blob \\
\lcf{5} & BEGIN OPENSSH PRIVATE KEY lines (one of them a real key) \\
\lcf{10} & OPENSSH PRIVATE KEY literals (BEGIN+END) \\
\lcf{8} & BEGIN PUBLIC KEY blocks (one is a literal with no body) \\
\lcf{8} & SPKI blocks that parse as keys \\
\lcf{6} & distinct RSA moduli \\
\lcf{429029376} & offset of the 2048-bit key \\
\lcf{740733952} & offset of the sshd\_config member in the tar table of contents \\
\lcf{3094} & size of sshd\_config, bytes \\
\midrule
\multicolumn{2}{l}{\lcf{BODYDATA.DAT}\quad\textbf{sha256 dbdd8b02ed81d6f0}} \\
\midrule
\lcf{272} & body tail: total bytes \\
\lcf{89a85276c720\allowbreak 8ea533a0e187\allowbreak fac4092c} & IV in the tail (in clear text) \\
\lcf{256} & signature block, bytes (= the size of an RSA-2048 signature) \\
\lcf{743818348} & offset of the tail in the file \\
\lcf{497f} & first bytes after raising the tail to the power e (no v1.5 structure) \\
\midrule
\multicolumn{2}{l}{\lcf{publication/\allowbreak live\_witness\allowbreak /FINGERPRINT\allowbreak \_RESULT\_2026\allowbreak -10-04\_1743.\allowbreak txt}\quad\textbf{sha256 411ec1397eb43f4f}} \\
\midrule
\lcf{192.168.1.10\allowbreak 2} & camera address in the runs \\
\lcf{MISMATCH} & verdict of the host-key comparison \\
\lcf{SHA256:0VOxJ\allowbreak qn75lys4Lf91\allowbreak tSlZu8nazWWx\allowbreak D6rTVc35jZo8\allowbreak 7U} & fingerprint of the camera live key \\
\lcf{d153b126a9fb\allowbreak e65cace0b7fd\allowbreak d6d4a566ef27\allowbreak 6b3596c43eab\allowbreak 4d5737e63668\allowbreak f3b5} & sha256 of the live key blob \\
\midrule
\multicolumn{2}{l}{\lcf{publication/\allowbreak live\_witness\allowbreak /live\_probe\_\allowbreak log.txt}\quad\textbf{sha256 4cbdc8c67c01054c}} \\
\midrule
\lcf{22} & probes that found port 22 open \\
\midrule
\multicolumn{2}{l}{\lcf{-}\quad\textbf{sha256 -}} \\
\midrule
\lcf{22} & camera sshd port \\
\midrule
\multicolumn{2}{l}{\lcf{publication/\allowbreak live\_witness\allowbreak /FINGERPRINT\allowbreak \_RESULT\_2026\allowbreak -10-04\_1743.\allowbreak txt}\quad\textbf{sha256 411ec1397eb43f4f}} \\
\midrule
\lcf{SHA256:J8L9a\allowbreak BLEimTdDW5m1\allowbreak iV7wU64LP1iW\allowbreak o3+lj9MGCxNO\allowbreak DI} & fingerprint of the firmware key it was compared against \\
\midrule
\multicolumn{2}{l}{\lcf{-}\quad\textbf{sha256 -}} \\
\midrule
\lcf{64} & of the 66 bytes of the key point differ (external measurement) \\
\lcf{66} & length of a P-256 point, bytes \\
\lcf{15740} & PTP/IP port (closed in the measured state) \\
\midrule
\multicolumn{2}{l}{\lcf{publication/\allowbreak verification\allowbreak \_runs/checks\allowbreak .txt}\quad\textbf{sha256 41efd26400c45698}} \\
\midrule
\lcf{15} & checks.py: green \\
\lcf{3} & checks.py: controls \\
\midrule
\multicolumn{2}{l}{\lcf{publication/\allowbreak verification\allowbreak \_runs/round2\allowbreak 77.txt}\quad\textbf{sha256 385c3674d0bff57f}} \\
\midrule
\lcf{13} & round277\_check.py: green \\
\lcf{5} & round277\_check.py: controls red \\
\midrule
\multicolumn{2}{l}{\lcf{publication/\allowbreak verification\allowbreak \_runs/round2\allowbreak 78.txt}\quad\textbf{sha256 78d1932d25d8911b}} \\
\midrule
\lcf{17} & round278\_check.py: green \\
\lcf{5} & round278\_check.py: controls red \\
\midrule
\multicolumn{2}{l}{\lcf{publication/\allowbreak verification\allowbreak \_runs/round2\allowbreak 79.txt}\quad\textbf{sha256 0a665fe918405af3}} \\
\midrule
\lcf{28} & round279\_check.py: green \\
\lcf{29} & round279\_check.py: total \\
\midrule
\multicolumn{2}{l}{\lcf{publication/\allowbreak verification\allowbreak \_runs/verify\allowbreak \_findings.tx\allowbreak t}\quad\textbf{sha256 30af8e4f5c111ad5}} \\
\midrule
\lcf{47} & verify\_findings.py: green \\
\lcf{5} & verify\_findings.py: controls red \\
\midrule
\multicolumn{2}{l}{\lcf{-}\quad\textbf{sha256 -}} \\
\midrule
\lcf{4} & number of corrections to numbers recorded earlier \\
\lcf{27} & instrument errors across the work (F140-F166) \\
\bottomrule
\end{longtable}
\endgroup


\end{document}
